\documentclass[a4paper,10pt]{article}

% Настройки русского языка и кодировок
\usepackage[T2A]{fontenc}
\usepackage[utf8]{inputenc}
\usepackage[russian]{babel}

% Математические пакеты
\usepackage{amsmath, amsfonts, amssymb, amsthm}
\usepackage{geometry}
\usepackage{titlesec}

% Настройка полей
\geometry{left=1cm, right=1cm, top=1.5cm, bottom=1.5cm}

% Настройка заголовков
\titleformat{\section}{\large\bfseries}{}{0em}{}
\titleformat{\subsection}{\normalsize\bfseries}{}{0em}{}

\begin{document}

\begin{center}
    \LARGE \textbf{Уравнения Пуассона и Лапласа на плоскости (2D)}
\end{center}

\section{Часть 1: Практическая (Алгоритмы и факты)}

\subsection*{Постановка задачи}
Рассматривается краевая задача для уравнения Лапласа ($\Delta u = 0$) или Пуассона ($\Delta u = f$) в двумерной области $G \subset \mathbb{R}^2$. Для полярной системы координат $(r, \varphi)$ оператор Лапласа имеет вид:
\[ \Delta u = u_{rr} + \frac{1}{r}u_r + \frac{1}{r^2}u_{\varphi\varphi} \]
Выделяют три основных типа областей и соответствующие им краевые задачи:
\begin{enumerate}
    \item \textbf{В круге} ($r < R$): Граничное условие ставится на окружности $r=R$. Неявно предполагается условие ограниченности решения в центре: $|u(0, \varphi)| < \infty$.
    \item \textbf{Во внешности круга} ($r > R$): Граничное условие на $r=R$. Неявно предполагается условие ограниченности решения на бесконечности: $|u(\infty, \varphi)| < \infty$.
    \item \textbf{В кольце} ($R_1 < r < R_2$): Задаются два граничных условия: на внутренней окружности $r=R_1$ и на внешней $r=R_2$.
\end{enumerate}

\textbf{Типы краевых условий} (задаются через производную по направлению внешней нормали $\vec{n}$):
\begin{itemize}
    \item \textit{Дирихле (1-й род):} задана сама функция $u|_{\partial G} = \mu(\varphi)$.
    \item \textit{Неймана (2-й род):} задана производная $u_n|_{\partial G} = \nu(\varphi)$.
    \item \textit{3-й род:} $(\alpha u + \beta u_n)|_{\partial G} = \mu(\varphi)$.
\end{itemize}
\textit{Внимание к знакам нормали:} Для круга внешняя нормаль сонаправлена с радиусом ($\partial_n = \partial_r$). Для внешности круга нормаль направлена к центру ($\partial_n = -\partial_r$). Для кольца на границе $R_1$ имеем $\partial_n = -\partial_r$, а на границе $R_2$ имеем $\partial_n = \partial_r$.

\subsection*{Алгоритм решения (Метод Фурье)}

\textbf{Шаг 0. Сведение уравнения Пуассона к уравнению Лапласа (если $\Delta u \neq 0$)} \\
Если задано уравнение Пуассона $\Delta u = f(r, \varphi)$, необходимо найти \textbf{любое} частное решение $w(r, \varphi)$ этого уравнения. Тогда делается замена $u(r, \varphi) = v(r, \varphi) + w(r, \varphi)$.
Функция $v(r, \varphi)$ будет удовлетворять однородному уравнению Лапласа $\Delta v = 0$ с модифицированными граничными условиями.

\textbf{Шаг 1. Разложение краевых условий в ряды Фурье} \\
Правые части граничных условий необходимо разложить в классические тригонометрические ряды Фурье по углу $\varphi$ на периоде $[0, 2\pi]$, являющемся базисом для угловой части решения (см. доказательство в Части 2):
\[ \mu(\varphi) = a_0 + \sum_{n=1}^\infty \left( a_n \cos n\varphi + b_n \sin n\varphi \right) \]

\textbf{Шаг 2. Выбор общего вида решения $\Delta u = 0$} \\
В зависимости от области $G$, решение ищется в виде линейной комбинации базисных гармонических функций. Коэффициенты, вызывающие сингулярности, зануляются (см. доказательство в Части 2):

\noindent
\textbf{А) В кольце ($R_1 < r < R_2$):} Участвуют все базисные функции.
\[ u(r, \varphi) = A_0 + B_0 \ln r + \sum_{n=1}^\infty \left[ \left( A_n r^n + B_n r^{-n} \right) \cos n\varphi + \left( C_n r^n + D_n r^{-n} \right) \sin n\varphi \right] \]
\textbf{Б) В круге ($r < R$):} Отбрасываются слагаемые, расходящиеся при $r \to 0$.
\[ u(r, \varphi) = A_0 + \sum_{n=1}^\infty r^n \big( A_n \cos n\varphi + C_n \sin n\varphi \big) \]
\textbf{В) Во внешности круга ($r > R$):} Отбрасываются слагаемые, расходящиеся при $r \to \infty$.
\[ u(r, \varphi) = A_0 + \sum_{n=1}^\infty r^{-n} \big( B_n \cos n\varphi + D_n \sin n\varphi \big) \]

\textbf{Шаг 3. Подстановка в граничные условия и поиск коэффициентов} \\
Выбранный ряд (и его производную по радиусу, если это нужно по условию) подставляют в граничные условия. Затем приравнивают коэффициенты при одинаковых гармониках ($\cos n\varphi$ и $\sin n\varphi$) в левой и правой частях уравнения.
\begin{itemize}
    \item Для кольца возникнут системы из двух уравнений на каждую пару констант (например, на $A_n$ и $B_n$).
    \item Для круга и внешности круга неизвестные константы выражаются явно из одного уравнения.
\end{itemize}
\textit{Замечание о разрешимости:} Если полученная алгебраическая система для коэффициентов имеет единственное решение — задача имеет единственное решение. Если хотя бы при одном $n$ система не имеет решений (часто возникает в задачах Неймана), то вся задача неразрешима. Если решений бесконечно много — задача имеет семейство решений.

\newpage
\section{Часть 2: Теоретическая (Обоснования и доказательства)}

\subsection*{Разделение переменных в полярных координатах}

\begin{proof}
Рассмотрим уравнение Лапласа $\Delta u = 0$ в полярных координатах:
\[ u_{rr} + \frac{1}{r}u_r + \frac{1}{r^2}u_{\varphi\varphi} = 0. \]
Будем искать нетривиальное решение в виде произведения двух функций: $u(r, \varphi) = R(r)\Phi(\varphi)$. Подставим в уравнение:
\[ R''(r)\Phi(\varphi) + \frac{1}{r}R'(r)\Phi(\varphi) + \frac{1}{r^2}R(r)\Phi''(\varphi) = 0. \]
Разделим обе части на $\frac{1}{r^2}R(r)\Phi(\varphi)$ и разнесем переменные:
\[ \frac{r^2 R''(r) + r R'(r)}{R(r)} = -\frac{\Phi''(\varphi)}{\Phi(\varphi)}. \]
Так как левая часть зависит только от $r$, а правая — только от $\varphi$, обе части должны быть равны некоторой константе разделения, которую обозначим как $\lambda$. Получаем два обыкновенных дифференциальных уравнения:
\begin{align}
    \Phi''(\varphi) + \lambda \Phi(\varphi) &= 0 \label{eq:phi} \\
    r^2 R''(r) + r R'(r) - \lambda R(r) &= 0 \label{eq:R}
\end{align}
\end{proof}

\subsection*{Уравнение по угловой координате (Периодичность)}

\begin{proof}
Функция $\Phi(\varphi)$ описывает распределение по полярному углу. Из физического (геометрического) смысла задачи следует, что классическое решение должно быть однозначным и гладким во всей плоскости. Это означает, что функция должна быть строго $2\pi$-периодической: $\Phi(\varphi + 2\pi) = \Phi(\varphi)$, а также непрерывно дифференцируемой (при «сшивке» краев $\varphi=0$ и $\varphi=2\pi$ должны совпадать значения функции и её производных).
Исследуем уравнение (1) при различных $\lambda$:
\begin{enumerate}
    \item \textbf{Случай $\lambda < 0$:} Обозначим $\omega = \sqrt{|\lambda|} > 0$. Общее решение имеет вид $\Phi(\varphi) = A e^{\omega \varphi} + B e^{-\omega \varphi}$. Эта функция является бесконечно гладкой на всей прямой. Чтобы она описывала гладкую $2\pi$-периодическую функцию, по теореме Ролля её производная $\Phi'(\varphi)$ должна обращаться в нуль как минимум в двух точках на периоде (в точках локального максимума и минимума). Однако уравнение $\Phi'(\varphi) = A\omega e^{\omega \varphi} - B\omega e^{-\omega \varphi} = 0$ может иметь не более одного действительного корня. Следовательно, нетривиальную гладкую периодическую функцию получить невозможно, и мы вынуждены положить $A = B = 0$.
    \item \textbf{Случай $\lambda = 0$:} Решение есть линейная функция $\Phi(\varphi) = A\varphi + B$. Требование совпадения значений на концах периода $\Phi(0) = \Phi(2\pi)$ дает: $B = A \cdot 2\pi + B \implies 2\pi A = 0 \implies A = 0$. Значит, периодическим решением здесь является только константа $\Phi_0(\varphi) = B$.
    \item \textbf{Случай $\lambda > 0$:} Обозначим $\lambda = \omega^2 > 0$. Общее решение имеет вид $\Phi(\varphi) = A \cos(\omega \varphi) + B \sin(\omega \varphi)$. Предположим, что решение нетривиально, то есть $A^2 + B^2 > 0$. Свернем эту линейную комбинацию в один косинус с фазой:
    \[ \Phi(\varphi) = \sqrt{A^2 + B^2} \left( \frac{A}{\sqrt{A^2 + B^2}} \cos(\omega\varphi) + \frac{B}{\sqrt{A^2 + B^2}} \sin(\omega\varphi) \right) = \sqrt{A^2 + B^2} \cos(\omega\varphi - \theta), \]
    где вспомогательный угол $\theta$ однозначно определяется из условий $\cos\theta = \frac{A}{\sqrt{A^2+B^2}}$ и $\sin\theta = \frac{B}{\sqrt{A^2+B^2}}$.
    Из условия $2\pi$-периодичности получаем тождество:
    \[ \cos(\omega\varphi - \theta) \equiv \cos\big(\omega(\varphi + 2\pi) - \theta\big). \]
    Поскольку это равенство должно выполняться при любом $\varphi$, подставим частное значение $\varphi = \frac{\theta}{\omega}$:
    \[ \cos(0) = \cos(2\pi\omega) \implies 1 = \cos(2\pi\omega). \]
    Отсюда строго следует, что $2\pi\omega = 2\pi n$, где $n \in \mathbb{N}$ (так как $\omega > 0$). Таким образом, $\omega = n$, что дает нам собственные значения $\lambda_n = n^2$, а собственные функции имеют классический вид $\Phi_n(\varphi) = A_n \cos(n\varphi) + B_n \sin(n\varphi)$.
\end{enumerate}
Следовательно, собственные значения равны $\lambda_n = n^2$ (где $n = 0, 1, 2, \dots$), а соответствующие им угловые функции образуют классическую тригонометрическую систему $\{1, \cos n\varphi, \sin n\varphi\}$.
\end{proof}

\textbf{Почему мы ищем решения только в таком виде?}
Функции $\{1, \cos n\varphi, \sin n\varphi\}$ суть собственные функции оператора второй производной на окружности (периодические краевые условия). По теореме Рисса-Фишера эта классическая тригонометрическая система образует \textbf{полный базис} в пространстве $L_2(S^1)$. Следовательно, любое гладкое граничное условие можно разложить по этому базису, и получающийся ряд полностью покрывает пространство всех возможных гармонических функций в круге/кольце. Аналогичные рассуждения проводились для начальной краевой задачи на отрезке для волнового уравнения.

\subsection*{Радиальное уравнение (Уравнение Эйлера)}

\begin{proof}
Подставим найденные значения $\lambda = n^2$ в радиальное уравнение \eqref{eq:R}:
\[ r^2 R''(r) + r R'(r) - n^2 R(r) = 0. \]
Это классическое уравнение Эйлера. Для его решения сделаем стандартную замену независимой переменной $r = e^t$, откуда $t = \ln r$. 
Вычислим производные сложной функции:
\[ R'(r) = \frac{dR}{dr} = \frac{dR}{dt} \frac{dt}{dr} = R'_t \cdot \frac{1}{r}, \]
\[ R''(r) = \frac{d}{dr}\left( R'_t \cdot \frac{1}{r} \right) = \left( \frac{d}{dt}(R'_t) \frac{dt}{dr} \right) \frac{1}{r} - R'_t \cdot \frac{1}{r^2} = R''_{tt} \cdot \frac{1}{r^2} - R'_t \cdot \frac{1}{r^2} = \frac{1}{r^2} \big( R''_{tt} - R'_t \big). \]
Подставим эти выражения в исходное уравнение:
\[ r^2 \left[ \frac{1}{r^2} \big( R''_{tt} - R'_t \big) \right] + r \left[ \frac{1}{r} R'_t \right] - n^2 R = 0 \implies R''_{tt} - R'_t + R'_t - n^2 R = 0. \]
Уравнение упрощается до линейного ОДУ второго порядка с постоянными коэффициентами:
\[ R''_{tt}(t) - n^2 R(t) = 0. \]
Характеристическое уравнение имеет вид $\mu^2 - n^2 = 0$. Рассмотрим два случая:
\begin{enumerate}
    \item \textbf{При $n \neq 0$:} Корни характеристического уравнения $\mu_{1,2} = \pm n$. 
    Общее решение имеет вид $R_n(t) = C_1 e^{nt} + C_2 e^{-nt}$. 
    Возвращаясь к переменной $r$ (с учетом $e^{nt} = (e^t)^n = r^n$), получаем:
    \[ R_n(r) = C_1 r^n + C_2 r^{-n}. \]
    \item \textbf{При $n = 0$:} Характеристическое уравнение имеет кратный корень $\mu_1 = \mu_2 = 0$. 
    В этом случае фундаментальная система решений состоит из константы и функции, линейной по $t$. Общее решение: $R_0(t) = C_1 + C_2 t$. 
    Сделав обратную замену $t = \ln r$, явно получаем логарифмическое слагаемое:
    \[ R_0(r) = C_1 + C_2 \ln r. \]
\end{enumerate}
Комбинируя радиальные решения $R_n(r)$ и угловые $\Phi_n(\varphi)$, мы получаем полную систему базисных гармонических функций на плоскости:
\[ \{1, \ln r, \ r^n \cos n\varphi, \ r^{-n} \cos n\varphi, \ r^n \sin n\varphi, \ r^{-n} \sin n\varphi \}. \]
\end{proof}

\subsection*{Обоснование отбрасывания сингулярных решений (Шаг 2)}

\begin{proof}
Физически осмысленное решение $u(r, \varphi)$ должно быть непрерывным и ограниченным внутри рассматриваемой области $G$.

1. \textbf{Для внутренности круга ($r < R$)} точка начала координат $r=0$ принадлежит области. 
Рассмотрим функции $\ln r$ и $r^{-n}$ (где $n \ge 1$). При $r \to 0$ обе эти функции стремятся к бесконечности: $\lim_{r \to 0} \ln r = -\infty$ и $\lim_{r \to 0} r^{-n} = \infty$. Чтобы решение оставалось ограниченным (классическим) в центре, коэффициенты при этих функциях обязаны быть равны нулю: $B_0 = 0$, $B_n = D_n = 0$.

2. \textbf{Для внешности круга ($r > R$)} область не ограничена, и координата $r \to \infty$ является частью области (точнее, границей на бесконечности). 
Рассмотрим функции $\ln r$ и $r^n$ (где $n \ge 1$). При $r \to \infty$ они неограниченно растут. Требование ограниченности решения на бесконечности $|u(\infty, \varphi)| < \infty$ заставляет нас положить коэффициенты при этих функциях равными нулю: $B_0 = 0$ (так как $\ln r \to \infty$), и $A_n = C_n = 0$ (так как $r^n \to \infty$). 

3. \textbf{Для кольца ($R_1 < r < R_2$)} ни точка $r=0$, ни $r=\infty$ не принадлежат области. Ни одна из базисных функций не обращается в бесконечность внутри кольца, поэтому ни один из коэффициентов заранее не зануляется, и необходимо искать все 4 набора коэффициентов $A_n, B_n, C_n, D_n$.
\end{proof}

\end{document}